\documentclass[final]{beamer}

% =========================================================
% POSTER FORMAT
% =========================================================
% A0 portrait:
\usepackage[
    orientation=portrait,
    size=a0,
    scale=1.00
]{beamerposter}

% =========================================================
% PACKAGES
% =========================================================
\usepackage[T1]{fontenc}
\usepackage[utf8]{inputenc}
\usepackage{lmodern}
\usepackage{graphicx}
\usepackage{booktabs}
\usepackage{array}
\usepackage{ragged2e}
\usepackage{tikz}
\usepackage{amsmath}
\usepackage{xcolor}
\usepackage{ifthen}

% =========================================================
% DISC LAB COLOUR PALETTE
% =========================================================
\definecolor{discnavy}{HTML}{1F2937}
\definecolor{discteal}{HTML}{2DD4BF}
\definecolor{discdarkteal}{HTML}{0F766E}
\definecolor{disclight}{HTML}{E5E7EB}
\definecolor{discsoft}{HTML}{F8FAFC}
\definecolor{discgray}{HTML}{6B7280}
\definecolor{discbluegray}{HTML}{334155}
\definecolor{ntucrimson}{HTML}{C8102E}

% =========================================================
% POSTER INFORMATION
% =========================================================
\title{Poster title}
\author{Author names}
\institute{%
Digital Intelligence for Structures and Construction Laboratory\\
School of Civil and Environmental Engineering\\
Nanyang Technological University, Singapore
}

% Replace these names with the actual logo filenames.
\newcommand{\disclogofile}{disc_lab_logo_white.pdf}
\newcommand{\ntulogofile}{ntu_logo_white.pdf}

% =========================================================
% TYPOGRAPHY
% =========================================================
\setbeamerfont{title}{
    series=\bfseries,
    size=\fontsize{48}{54}
}

\setbeamerfont{author}{
    series=\bfseries,
    size=\fontsize{24}{29}
}

\setbeamerfont{institute}{
    size=\fontsize{18}{23}
}

\setbeamerfont{block title}{
    series=\bfseries,
    size=\fontsize{25}{30}
}

\setbeamerfont{block body}{
    size=\fontsize{18}{23}
}

\setbeamerfont{caption}{
    size=\fontsize{15}{19}
}

% =========================================================
% BEAMER COLOURS
% =========================================================
\setbeamercolor{background canvas}{
    bg=white
}

\setbeamercolor{normal text}{
    fg=discnavy,
    bg=white
}

\setbeamercolor{headline}{
    fg=white,
    bg=discnavy
}

\setbeamercolor{footline}{
    fg=white,
    bg=discnavy
}

\setbeamercolor{block title}{
    fg=white,
    bg=discnavy
}

\setbeamercolor{block body}{
    fg=discnavy,
    bg=discsoft
}

\setbeamercolor{block title alerted}{
    fg=discnavy,
    bg=discteal
}

\setbeamercolor{block body alerted}{
    fg=discnavy,
    bg=white
}

\setbeamercolor{block title example}{
    fg=white,
    bg=discdarkteal
}

\setbeamercolor{block body example}{
    fg=discnavy,
    bg=white
}

% =========================================================
% GENERAL FORMATTING
% =========================================================
\setbeamertemplate{navigation symbols}{}
\setbeamertemplate{blocks}[rounded][shadow=false]

\setbeamersize{
    text margin left=1.4cm,
    text margin right=1.4cm
}

\setlength{\parskip}{0.35em}
\setlength{\columnsep}{1.0cm}

\setbeamertemplate{itemize item}{
    \color{discteal}\raisebox{0.20ex}{\Large$\blacktriangleright$}
}

\setbeamertemplate{itemize subitem}{
    \color{discdarkteal}\raisebox{0.15ex}{\large$\bullet$}
}

\setbeamertemplate{enumerate item}{
    \color{discdarkteal}\insertenumlabel.
}

% =========================================================
% LOGO COMMAND
% =========================================================
\newcommand{\posterlogo}[2]{%
    \IfFileExists{#1}{%
        \includegraphics[
            width=#2,
            height=6.5cm,
            keepaspectratio
        ]{#1}
    }{%
        \fcolorbox{white}{discnavy}{%
            \begin{minipage}[c][5.4cm][c]{#2}
                \centering
                \color{white}
                \bfseries
                Logo placeholder
            \end{minipage}
        }
    }%
}

% =========================================================
% IMAGE PLACEHOLDER
% =========================================================
\newcommand{\imgplaceholder}[2]{%
    \fcolorbox{discgray}{white}{%
        \begin{minipage}[c][#1][c]{0.96\linewidth}
            \centering

            {\fontsize{23}{27}\selectfont
            \color{discgray}
            \bfseries
            IMAGE PLACEHOLDER}

            \vspace{0.35cm}

            {\fontsize{16}{20}\selectfont
            \color{discgray}
            #2}
        \end{minipage}
    }
}

% =========================================================
% FIGURE COMMAND
% =========================================================
% Usage:
% \posterfigure{filename}{height}{caption}
%
% The placeholder appears automatically when the file is absent.
\newcommand{\posterfigure}[3]{%
    \begin{center}
        \IfFileExists{#1}{%
            \includegraphics[
                width=\linewidth,
                height=#2,
                keepaspectratio
            ]{#1}
        }{%
            \imgplaceholder{#2}{Expected file: \texttt{#1}}
        }

        \vspace{0.20cm}

        \begin{minipage}{0.96\linewidth}
            \centering
            \fontsize{15}{19}\selectfont
            \color{discgray}
            #3
        \end{minipage}
    \end{center}
}

% =========================================================
% SECTION BANNER
% =========================================================
\newcommand{\sectionbanner}[1]{%
    \begin{beamercolorbox}[
        wd=\linewidth,
        sep=0.55cm,
        rounded=true
    ]{block title}
        \centering
        \fontsize{27}{32}\selectfont
        \bfseries
        #1
    \end{beamercolorbox}
}

% =========================================================
% HEADER
% =========================================================
\setbeamertemplate{headline}{%
    \begin{beamercolorbox}[
        wd=\paperwidth,
        ht=10.8cm,
        dp=0cm
    ]{headline}

        \vspace{0.9cm}

        \begin{columns}[
            T,
            totalwidth=0.95\paperwidth
        ]

            \begin{column}{0.20\paperwidth}
                \centering
                \posterlogo{\disclogofile}{20cm}
            \end{column}

            \begin{column}{0.55\paperwidth}
                \centering

                {\usebeamerfont{title}
                \color{white}
                \inserttitle\par}

                \vspace{0.55cm}

                {\usebeamerfont{author}
                \color{discteal}
                \insertauthor\par}

                \vspace{0.35cm}

                {\usebeamerfont{institute}
                \color{white}
                \insertinstitute\par}
            \end{column}

            \begin{column}{0.20\paperwidth}
                \centering
                \posterlogo{\ntulogofile}{20.5cm}
            \end{column}

        \end{columns}

        \vspace{0.7cm}

        \begin{beamercolorbox}[
            wd=\paperwidth,
            ht=0.28cm,
            dp=0cm
        ]{block title alerted}
        \end{beamercolorbox}

    \end{beamercolorbox}
}

% =========================================================
% FOOTER
% =========================================================
\setbeamertemplate{footline}{%
    \begin{beamercolorbox}[
        wd=\paperwidth,
        ht=2.2cm,
        dp=0.7cm,
        leftskip=1.5cm,
        rightskip=1.5cm
    ]{footline}

        \vspace{0.45cm}

        \fontsize{16}{20}\selectfont
        \textbf{DISC Lab}
        \hfill
        \texttt{disclab.pantoja-rosero.com}
        \hfill
        \texttt{bryan.pr@ntu.edu.sg}

    \end{beamercolorbox}
}

% =========================================================
% DOCUMENT
% =========================================================
\begin{document}

\begin{frame}[t]

\vspace{0.65cm}

% ---------------------------------------------------------
% FULL-WIDTH OPENING SECTION
% ---------------------------------------------------------
\begin{alertblock}{Research context and central message}
    \justifying

    Introduce the engineering problem and explain why the research is
    relevant. This section should normally contain two or three concise
    sentences.

    State the central contribution of the research clearly. A reader
    should understand the main message of the poster without reading
    the remaining sections.
\end{alertblock}

\vspace{0.45cm}

% ---------------------------------------------------------
% FULL-WIDTH OVERVIEW FIGURE
% ---------------------------------------------------------
\begin{block}{Research framework}
    \justifying

    Briefly introduce the common methodological framework or the
    relationship among the research components.

    \posterfigure
        {figures/research_framework.png}
        {14cm}
        {Overview of the research framework and its principal outputs.}
\end{block}

\vspace{0.45cm}

% ---------------------------------------------------------
% TWO-COLUMN VERTICAL BODY
% ---------------------------------------------------------
\begin{columns}[t,totalwidth=\textwidth]

    % =====================================================
    % LEFT COLUMN
    % =====================================================
    \begin{column}{0.485\textwidth}

        \begin{block}{Research component 1}
            \justifying

            Introduce the first research component in two or three
            sentences. Emphasise its input, methodology, output, and
            intended engineering use.

            \posterfigure
                {figures/component_01.png}
                {13cm}
                {Suggested caption for research component 1.}
        \end{block}

        \begin{block}{Research component 2}
            \justifying

            Introduce the second research component. Keep the text
            focused on the specific engineering problem addressed and
            the technological contribution.

            \posterfigure
                {figures/component_02.png}
                {13cm}
                {Suggested caption for research component 2.}
        \end{block}

        \begin{block}{Research component 3}
            \justifying

            Introduce the third research component and explain how it
            extends or complements the previous work.

            \posterfigure
                {figures/component_03.png}
                {13cm}
                {Suggested caption for research component 3.}
        \end{block}

    \end{column}

    % =====================================================
    % RIGHT COLUMN
    % =====================================================
    \begin{column}{0.485\textwidth}

        \begin{block}{Research component 4}
            \justifying

            Introduce the fourth research component in two or three
            sentences. Highlight the interaction between digital data
            and engineering practice.

            \posterfigure
                {figures/component_04.png}
                {13cm}
                {Suggested caption for research component 4.}
        \end{block}

        \begin{block}{Research component 5}
            \justifying

            Introduce the fifth research component. Describe its
            operational or field-level relevance.

            \posterfigure
                {figures/component_05.png}
                {13cm}
                {Suggested caption for research component 5.}
        \end{block}

        \begin{block}{Research component 6}
            \justifying

            Introduce the sixth research component and connect it to
            engineering analysis, decision support, or automation.

            \posterfigure
                {figures/component_06.png}
                {13cm}
                {Suggested caption for research component 6.}
        \end{block}

    \end{column}

\end{columns}

\vspace{0.45cm}

% ---------------------------------------------------------
% FULL-WIDTH IMPACT SECTION
% ---------------------------------------------------------
\begin{exampleblock}{Engineering applications and impact}

    \begin{columns}[t,totalwidth=\linewidth]

        \begin{column}{0.24\linewidth}
            \centering
            \textbf{Inspection}

            \vspace{0.2cm}

            Objective and traceable
            condition assessment.
        \end{column}

        \begin{column}{0.24\linewidth}
            \centering
            \textbf{Asset management}

            \vspace{0.2cm}

            Structured digital records
            of existing assets.
        \end{column}

        \begin{column}{0.24\linewidth}
            \centering
            \textbf{Maintenance}

            \vspace{0.2cm}

            Better documentation and
            intervention planning.
        \end{column}

        \begin{column}{0.24\linewidth}
            \centering
            \textbf{Engineering analysis}

            \vspace{0.2cm}

            Digital information linked
            to computational models.
        \end{column}

    \end{columns}

\end{exampleblock}

\vspace{0.45cm}

% ---------------------------------------------------------
% CONCLUSION AND OUTPUTS
% ---------------------------------------------------------
\begin{columns}[t,totalwidth=\textwidth]

    \begin{column}{0.56\textwidth}

        \begin{alertblock}{Conclusion}
            \justifying

            Summarise the principal contribution in two or three
            sentences. State how the combined research advances
            engineering practice and identify the next stage of the
            research programme.
        \end{alertblock}

    \end{column}

    \begin{column}{0.42\textwidth}

        \begin{block}{Selected outputs}
            \justifying

            \begin{itemize}
                \item Selected journal publication.
                \item Selected conference publication.
                \item Dataset, software, or demonstrator.
            \end{itemize}
        \end{block}

    \end{column}

\end{columns}

\vspace{0.35cm}

% ---------------------------------------------------------
% ACKNOWLEDGEMENT
% ---------------------------------------------------------
\begin{block}{Acknowledgements and contact}
    \justifying

    Add funding agencies, collaborating institutions, industry
    partners, project identifiers, a QR code, and relevant contact
    information.
\end{block}

\end{frame}

\end{document}